\section{Introduction}
Positron emission tomography (PET) detects pairs of annihilation photons to reconstruct
radiotracer distributions. Time-of-flight (TOF) information uses the difference between
photon arrival times to constrain the annihilation position along the line joining the
detectors and can improve image signal-to-noise performance. The precision of this
information depends on the coincidence timing resolution (CTR), conventionally expressed
as the full width at half maximum (FWHM) of the time-difference distribution~\cite{surti2015}.

A leading-edge discriminator (LED) assigns an arrival time when a pulse crosses a fixed
amplitude threshold. Variations in pulse amplitude and rising-edge shape shift the crossing
time, producing time walk. Digitized waveforms retain information about these variations,
which machine learning (ML) can exploit to estimate event-specific timing corrections.

Waveform-based approaches include direct TOF estimation and correction of conventional
timing measurements. Berg and Cherry demonstrated direct TOF estimation from detector
waveform pairs using a convolutional neural network (CNN), improving CTR by 20\% relative
to LED, from 231 to 185 ps~\cite{berg2018}. Onishi et al. combined LED with a CNN
trained using waveforms acquired at a single source position, improving timing
resolution while avoiding the position-dependent bias associated with direct TOF
estimation~\cite{onishi2022}. Feng et al. explored convolutional features combined with
transformer-based attention to capture local and global temporal information~\cite{feng2024}.
For bismuth germanate (BGO) detectors, Loignon-Houle et al. compared LED, a
two-threshold time-walk correction and CNN-based timing, showing that performance gains
and timing-distribution tails depend on detector characteristics~\cite{loignon2025}.

An open question is whether waveform samples beyond the rising edge improve
ML-based LED timing corrections, and which temporal regions contribute most.
Previous CNN studies have primarily examined short windows near the rising
edge~\cite{onishi2022,loignon2025}, leaving extended-window performance at high
sampling rates insufficiently characterized. In related model-based work,
Ruiz-Gonzalez et al. observed diminishing returns from additional samples
and found that the advantage of maximum-likelihood timing over digital
discrimination disappeared above 2 GS/s~\cite{ruiz2018}. These results do
not establish whether extending the waveform window benefits data-driven
corrections at the 80 GS/s acquisition rate used here.

The analog readout chain is another determinant of timing performance. Amplification can
change signal amplitude, bandwidth, noise and rise time, affecting both LED timing
and the waveform features accessible to ML. Pourashraf et al. compared front-end
readout configurations for silicon photomultiplier (SiPM)-based TOF-PET detectors,
reporting a best CTR of $99.4 \pm 1.9$ ps, more than 20 ps better than two alternative
configurations~\cite{pourashraf2022}. These results establish the importance of
electronics choices, but do not determine whether ML corrections applied to outputs
of successive stages of cascaded amplification perform differently.

Building on previous waveform-based timing studies, this work investigates how different 
ML approaches exploit detector signals to correct leading-edge timing, rather than focusing 
exclusively on achieving the lowest CTR. 
By comparing multiple estimator families and temporal windows, we aim to characterize the timing 
information available across the waveform and assess the potential of both conventional and 
dedicated high-gain readouts.

A further objective is to investigate the scalability of ML-based timing calibration. 
Training a separate model for every detector pair would become impractical in complex PET systems. 
We therefore introduce detector-shared estimators, which apply the same correction function to both 
detector waveforms, enforcing antisymmetry under detector exchange. Comparing these models with 
unconstrained estimators allows us to assess whether detector-specific relationships are necessary 
or whether general waveform characteristics are sufficient for effective timing correction.

Although complete waveform acquisition may not be practical in clinical systems, this analysis 
provides a way to investigate which signal characteristics ML models exploit and which temporal 
regions are most relevant to their predictions. The two electronics boards are evaluated 
independently using control, development, and blind acquisitions, with temporal sensitivity 
analysis complementing the timing-performance evaluation.